% front_matter_definitions.tex -- Book I. Definitions (35 definitions, jemmybutton-matched).
% Figure geometry and text from jemmybutton byrne-en-latex.tex (CC-BY-SA).
% Build: lualatex front_matter_definitions.tex  ->  front_matter_definitions.pdf
\documentclass[12pt]{article}
\usepackage{geometry}
\geometry{a5paper, margin=1.5cm}
\usepackage[lining]{ebgaramond}
\usepackage{amssymb}
\usepackage{byrne}
\def\mpPre{u := 1cm; textLabels := false;}

% Helper: definition block without figure (full-width, centered)
\newcommand{\deftext}[2]{%
  \vspace{0.35cm}%
  \begin{center}%
    \textbf{#1.}\\[2pt]%
    #2%
  \end{center}%
}

% Helper: definition block with figure (figure left, text right)
\newcommand{\deffig}[2]{%
  \vspace{0.35cm}%
  \noindent%
  \begin{minipage}[t]{0.34\linewidth}%
    \vspace{0pt}\centering%
    \drawCurrentPicture%
  \end{minipage}\hfill%
  \begin{minipage}[t]{0.62\linewidth}%
    \vspace{0pt}\centering%
    \textbf{#1.}\\[4pt]%
    #2%
  \end{minipage}%
  \vspace{0.1cm}%
}

\begin{document}

\begin{center}
  {\large\textbf{Book~I.}}\\[4pt]
  {\Large\textbf{Definitions}}
\end{center}

\vspace{0.4cm}
\noindent\rule{\linewidth}{0.4pt}
\vspace{0.2cm}

% --- Defs 1–8: no figures ---

\deftext{1}{A \emph{point} is that which has no parts.}

\deftext{2}{A \emph{line} is length without breadth.}

\deftext{3}{The extremities of a line are points.}

\deftext{4}{The straight or right line is that which lies evenly between its extremities.}

\deftext{5}{A \emph{surface} is that which has length and breadth only.}

\deftext{6}{The extremities of a surface are lines.}

\deftext{7}{A \emph{plane surface} is that which lies evenly between its extremities.}

\deftext{8}{A \emph{plane angle} is the inclination of two lines to one another,
in a plane, which meet together, but are not in the same direction.}

% --- Def 9: plane rectilinear angle ---
\defineNewPicture{%
  pair A, B, C;%
  A := (0, 0);%
  B := (2/3u, 2/3u);%
  C := (u, ypart(A));%
  byAngleDefine(B, A, C, byyellow, SOLID_SECTOR);%
  draw byNamedAngleResized();%
  byLineDefine(A, B, byblue, SOLID_LINE, REGULAR_WIDTH);%
  byLineDefine(A, C, byred,  SOLID_LINE, REGULAR_WIDTH);%
  draw byNamedLineSeq(0)(AC,AB);%
}
\deffig{9}{A \emph{plane rectilinear angle} is the inclination of two straight lines
to one another, which meet together, but are not in the same straight line.}

% --- Def 10: right angle ---
\defineNewPicture{%
  pair A, B, C, D;%
  A := (0, 0);%
  B := (u, 0);%
  C := (0, 2/3u);%
  D := (-u, 0);%
  byAngleDefine(B, A, C, byblack, ARC_SECTOR);%
  byAngleDefine(D, A, C, byblack, ARC_SECTOR);%
  draw byNamedAngleResized();%
  draw byLine(D, B, byblack, SOLID_LINE, REGULAR_WIDTH);%
  draw byLine(A, C, byblack, SOLID_LINE, REGULAR_WIDTH);%
}
\deffig{10}{When one straight line standing on another straight line makes the
adjacent angles equal, each of these angles is called a \emph{right angle},
and each of these lines is said to be \emph{perpendicular} to one another.}

% --- Def 11: obtuse angle ---
\defineNewPicture{%
  pair A, B, C;%
  A := (0, 0);%
  B := (-2/3u, 2/3u);%
  C := (u, ypart(A));%
  byAngleDefine(B, A, C, byred, SOLID_SECTOR);%
  draw byNamedAngleResized();%
  byLineDefine(A, B, byyellow, SOLID_LINE, REGULAR_WIDTH);%
  byLineDefine(A, C, byblue,   SOLID_LINE, REGULAR_WIDTH);%
  draw byNamedLineSeq(0)(AC,AB);%
}
\deffig{11}{An \emph{obtuse angle} is an angle greater than a right angle.}

% --- Def 12: acute angle ---
\defineNewPicture{%
  pair A, B, C;%
  A := (0, 0);%
  B := (2/3u, 2/3u);%
  C := (u, ypart(A));%
  byAngleDefine(B, A, C, byblue, SOLID_SECTOR);%
  draw byNamedAngleResized();%
  byLineDefine(A, B, byyellow, SOLID_LINE, REGULAR_WIDTH);%
  byLineDefine(A, C, byred,    SOLID_LINE, REGULAR_WIDTH);%
  draw byNamedLineSeq(0)(AC,AB);%
}
\deffig{12}{An \emph{acute angle} is an angle less than a right angle.}

% --- Defs 13–14: no figures ---

\deftext{13}{A \emph{term} or \emph{boundary} is the extremity of any thing.}

\deftext{14}{A \emph{figure} is a surface enclosed on all sides by a line or lines.}

% --- Def 15: circle ---
\defineNewPicture{%
  pair O, A, B, C, D, E;%
  numeric r;%
  r := 1/2u;%
  O := (0, 0);%
  A := dir(0)   scaled r;%
  B := dir(60)  scaled r;%
  C := dir(130) scaled r;%
  D := dir(180) scaled r;%
  E := dir(-60) scaled r;%
  draw byLine(O, B)(byblack, SOLID_LINE, REGULAR_WIDTH);%
  draw byLine(O, C)(byred,   SOLID_LINE, REGULAR_WIDTH);%
  draw byLine(O, E)(byyellow,SOLID_LINE, REGULAR_WIDTH);%
  draw byLine(D, A)(byblue,  SOLID_LINE, REGULAR_WIDTH);%
  draw byCircleR(O, r, byred, 0, 0, 0);%
}
\deffig{15}{A \emph{circle} is a plane figure, bounded by one continued line,
called its circumference or periphery; and having a certain point within it,
from which all straight lines drawn to its circumference are equal.}

% --- Def 16: no figure ---
\deftext{16}{This point (from which the equal lines are drawn) is called the \emph{centre} of the circle.}

% --- Def 17: diameter ---
\defineNewPicture{%
  pair O, A, B;%
  numeric r;%
  r := 1/2u;%
  O := (0, 0);%
  A := dir(0)   scaled r;%
  B := dir(180) scaled r;%
  draw byLine(A, B)(byyellow, SOLID_LINE, REGULAR_WIDTH);%
  draw byCircleR(O, r, byred, 0, 0, 0);%
}
\deffig{17}{A \emph{diameter} of a circle is a straight line drawn through the
centre, terminated both ways in the circumference.}

% --- Def 18: semicircle ---
\defineNewPicture{%
  pair O, A, B;%
  numeric r;%
  r := 1/2u;%
  O := (0, 0);%
  A := dir(0)   scaled r;%
  B := dir(180) scaled r;%
  draw byLine(A, B)(byblue, SOLID_LINE, REGULAR_WIDTH);%
  draw byArc(O, A, B)(r, byyellow, 0, 0, 0, 0);%
  draw byArc(O, B, A)(r, byyellow, 1, 0, 0, 0);%
}
\deffig{18}{A \emph{semicircle} is the figure contained by the diameter, and
the part of the circle cut off by the diameter.}

% --- Def 19: segment ---
\defineNewPicture{%
  pair O, A, B;%
  path P;%
  numeric r;%
  r := 1/2u;%
  P := fullcircle scaled 2r;%
  O := (0, 0);%
  A := point 1 of P;%
  B := point 3 of P;%
  draw byLine(A, B)(byred, SOLID_LINE, REGULAR_WIDTH);%
  draw byArc(O, A, B)(r, byblue, 0, 0, 0, 0);%
  draw byArc(O, B, A)(r, byblue, 1, 0, 0, 0);%
}
\deffig{19}{A \emph{segment} of a circle is a figure contained by a straight
line and the part of the circumference which it cuts off.}

% --- Defs 20–21: no figures ---
\deftext{20}{A figure contained by straight lines only, is called a \emph{rectilinear figure}.}

\deftext{21}{A \emph{triangle} is a rectilinear figure included by three sides.}

% --- Def 22: quadrilateral (with figure) ---
\defineNewPicture{%
  pair A, B, C, D;%
  A := (0, 0);%
  B := (u, 1/2u);%
  C := (-1/2u, -4/3u);%
  D := (4/3u,  -u);%
  draw byLine(C, B)(byred,    SOLID_LINE, REGULAR_WIDTH);%
  draw byLine(A, D)(byblue,   SOLID_LINE, REGULAR_WIDTH);%
  byLineDefine(A, B, byyellow, SOLID_LINE, REGULAR_WIDTH);%
  byLineDefine(A, C, byyellow, SOLID_LINE, REGULAR_WIDTH);%
  byLineDefine(B, D, byyellow, SOLID_LINE, REGULAR_WIDTH);%
  byLineDefine(C, D, byblack,  SOLID_LINE, REGULAR_WIDTH);%
  draw byNamedLineSeq(0)(BD,CD,AC,AB);%
  draw byLabelsOnPolygon(A, B, D, C)(ALL_LABELS, 0);%
}
\deffig{22}{A \emph{quadrilateral figure} is one which is bounded by four sides.
The straight lines \drawUnitLine{AD} and \drawUnitLine{CB} connecting the
vertices of the opposite angles of a quadrilateral figure, are called its
\emph{diagonals}.}

% --- Def 23: no figure ---
\deftext{23}{A \emph{polygon} is a rectilinear figure bounded by more than four sides.}

% --- Def 24: equilateral triangle ---
\defineNewPicture{%
  pair A, B, C;%
  A := dir(-30)  scaled 1/2u;%
  B := dir(-150) scaled 1/2u;%
  C := dir(90)   scaled 1/2u;%
  byLineDefine(A, B, byblue,   SOLID_LINE, REGULAR_WIDTH);%
  byLineDefine(B, C, byred,    SOLID_LINE, REGULAR_WIDTH);%
  byLineDefine(C, A, byyellow, SOLID_LINE, REGULAR_WIDTH);%
  draw byNamedLineSeq(0)(AB,BC,CA);%
}
\deffig{24}{A triangle whose sides are equal, is said to be \emph{equilateral}.}

% --- Def 25: isosceles triangle ---
\defineNewPicture{%
  pair A, B, C;%
  A := dir(-60)  scaled 1/2u;%
  B := dir(-120) scaled 1/2u;%
  C := dir(90)   scaled 1/2u;%
  byLineDefine(A, B, byblue, SOLID_LINE, REGULAR_WIDTH);%
  byLineDefine(B, C, byred,  SOLID_LINE, REGULAR_WIDTH);%
  byLineDefine(C, A, byred,  SOLID_LINE, REGULAR_WIDTH);%
  draw byNamedLineSeq(0)(AB,BC,CA);%
}
\deffig{25}{A triangle which has only two sides equal is called an \emph{isosceles triangle}.}

% --- Def 26: no figure ---
\deftext{26}{A \emph{scalene triangle} is one which has no two sides equal.}

% --- Def 27: right-angled triangle ---
\defineNewPicture{%
  pair A, B, C;%
  A := (0, 0);%
  B := (-u, 0);%
  C := (0, 3/4u);%
  byLineDefine(A, B, byred,    SOLID_LINE, REGULAR_WIDTH);%
  byLineDefine(B, C, byyellow, SOLID_LINE, REGULAR_WIDTH);%
  byLineDefine(C, A, byblue,   SOLID_LINE, REGULAR_WIDTH);%
  draw byNamedLineSeq(0)(AB,BC,CA);%
}
\deffig{27}{A \emph{right angled triangle} is that which has a right angle.}

% --- Def 28: obtuse-angled triangle ---
\defineNewPicture{%
  pair A, B, C;%
  A := (-1/4u, 0);%
  B := (-u, 0);%
  C := (0, 3/4u);%
  byLineDefine(A, B, byred,    SOLID_LINE, REGULAR_WIDTH);%
  byLineDefine(B, C, byblue,   SOLID_LINE, REGULAR_WIDTH);%
  byLineDefine(C, A, byyellow, SOLID_LINE, REGULAR_WIDTH);%
  draw byNamedLineSeq(0)(AB,BC,CA);%
}
\deffig{28}{An \emph{obtuse angled triangle} is that which has an obtuse angle.}

% --- Def 29: acute-angled triangle ---
\defineNewPicture{%
  pair A, B, C;%
  A := (0, 0);%
  B := (-u, 0);%
  C := (-1/4u, 3/4u);%
  byLineDefine(A, B, byblue,   SOLID_LINE, REGULAR_WIDTH);%
  byLineDefine(B, C, byyellow, SOLID_LINE, REGULAR_WIDTH);%
  byLineDefine(C, A, byred,    SOLID_LINE, REGULAR_WIDTH);%
  draw byNamedLineSeq(0)(AB,BC,CA);%
}
\deffig{29}{An \emph{acute angled triangle} is that which has three acute angles.}

% --- Def 30: square ---
\defineNewPicture{%
  pair A, B, C, D;%
  numeric s;%
  s := u;%
  A := (0, 0);%
  B := (s, 0);%
  C := (0, s);%
  D := (s, s);%
  byLineDefine(A, B, byred,    SOLID_LINE, REGULAR_WIDTH);%
  byLineDefine(A, C, byblue,   SOLID_LINE, REGULAR_WIDTH);%
  byLineDefine(B, D, byyellow, SOLID_LINE, REGULAR_WIDTH);%
  byLineDefine(C, D, byblack,  SOLID_LINE, REGULAR_WIDTH);%
  draw byNamedLineSeq(0)(AB,AC,CD,BD);%
}
\deffig{30}{Of four-sided figures, a \emph{square} is that which has all its sides equal,
and all its angles right angles.}

% --- Def 31: rhombus ---
\defineNewPicture{%
  pair A, B, C, D;%
  numeric s;%
  s := u;%
  A := (0, 0);%
  B := (s, 0);%
  C := A shifted (dir(80) scaled s);%
  D := B shifted (dir(80) scaled s);%
  byLineDefine(A, B, byred,    SOLID_LINE, REGULAR_WIDTH);%
  byLineDefine(A, C, byblue,   SOLID_LINE, REGULAR_WIDTH);%
  byLineDefine(B, D, byyellow, SOLID_LINE, REGULAR_WIDTH);%
  byLineDefine(C, D, byblack,  SOLID_LINE, REGULAR_WIDTH);%
  draw byNamedLineSeq(0)(AB,AC,CD,BD);%
}
\deffig{31}{A \emph{rhombus} is that which has all its sides equal, but its angles
are not right angles.}

% --- Def 32: oblong ---
\defineNewPicture{%
  pair A, B, C, D;%
  numeric s;%
  s := u;%
  A := (0, 0);%
  B := (4/3s, 0);%
  C := (0, 3/4s);%
  D := (4/3s, 3/4s);%
  byLineDefine(A, B, byblue,  SOLID_LINE, REGULAR_WIDTH);%
  byLineDefine(A, C, byred,   SOLID_LINE, REGULAR_WIDTH);%
  byLineDefine(B, D, byred,   SOLID_LINE, REGULAR_WIDTH);%
  byLineDefine(C, D, byblue,  SOLID_LINE, REGULAR_WIDTH);%
  draw byNamedLineSeq(0)(AB,AC,CD,BD);%
}
\deffig{32}{An \emph{oblong} is that which has all its angles right angles, but
has not all its sides equal.}

% --- Def 33: rhomboid ---
\defineNewPicture{%
  pair A, B, C, D;%
  numeric s;%
  s := u;%
  A := (0, 0);%
  B := (s, 0);%
  C := (1/4s, 3/4s);%
  D := (s + 1/4s, 3/4s);%
  byLineDefine(A, B, byblue, SOLID_LINE, REGULAR_WIDTH);%
  byLineDefine(A, C, byred,  SOLID_LINE, REGULAR_WIDTH);%
  byLineDefine(B, D, byred,  SOLID_LINE, REGULAR_WIDTH);%
  byLineDefine(C, D, byblue, SOLID_LINE, REGULAR_WIDTH);%
  draw byNamedLineSeq(0)(AB,AC,CD,BD);%
}
\deffig{33}{A \emph{rhomboid} is that which has its opposite sides equal to one
another, but all its sides are not equal nor its angles right angles.}

% --- Def 34: no figure ---
\deftext{34}{All other quadrilateral figures are called \emph{trapeziums}.}

% --- Def 35: parallel lines ---
\defineNewPicture{%
  pair A, B, C, D;%
  numeric s;%
  s := u;%
  A := (0, 0);%
  B := (4/3s, 0);%
  C := (0, 1/2s);%
  D := (4/3s, 1/2s);%
  draw byLine(A, B, byred,    SOLID_LINE, REGULAR_WIDTH);%
  draw byLine(C, D, byyellow, SOLID_LINE, REGULAR_WIDTH);%
}
\deffig{35}{\emph{Parallel} straight lines are such as are in the same plane,
and which being produced continually in both directions would never meet.}

\end{document}
